% =============================================================
%  Chapter — HVDC Topologies
%  Assets (SVG source + PDF build) live in this same subfolder.
% =============================================================
%\sectionslide{System Overview}{}



\begin{frame}{1GW - Symmetrical Monopole (SMP)}
  \begin{columns}[c,onlytextwidth]
    \begin{column}{0.46\textwidth}
      \centering
      \topoimg[width=\linewidth]{projects/fervo/smp}
    \end{column}
    \begin{column}{0.50\textwidth}
      \vspace{0.5em}
      \begin{itemize}
        \setlength{\itemsep}{0.7em}
        \item \textbf{Topology} - SMP with $\pm400kV$
        \item \textbf{DC Line} - Two symmetrical conductors
      \end{itemize}
    \end{column}
  \end{columns}
\end{frame}


\begin{frame}{3GW - Dual Symmetrical Monopole (DSMP)}
  \begin{columns}[c,onlytextwidth]
    \begin{column}{0.46\textwidth}
      \centering
      \topoimg[width=\linewidth]{projects/fervo/dsmp}
    \end{column}
    \begin{column}{0.50\textwidth}
      \vspace{0.5em}
      \begin{itemize}
        \setlength{\itemsep}{0.7em}
        \item \textbf{Topology} - DSMP with $\pm400kV$
        \item \textbf{DC Line} - Four symmetrical conductors
      \end{itemize}
    \end{column}
  \end{columns}
\end{frame}


\begin{frame}{3GW - Bipole (BIP)}
  \begin{columns}[c,onlytextwidth]
    \begin{column}{0.46\textwidth}
      \centering
      \topoimg[width=\linewidth]{projects/fervo/bip}
    \end{column}
    \begin{column}{0.50\textwidth}
      \vspace{0.5em}
      \begin{itemize}
        \setlength{\itemsep}{0.7em}
        \item \textbf{Topology} - Bipole with $\pm525kV$
        \item \textbf{DC Line}
        \begin{itemize}
          \item Two high voltage conductors (HV) and one dedicated metallic return conductor (DMR).
          \item In case of cable applications, parallel cables are required due to high current rating.
          \item The voltage rating of the DMR depends on the conductor resign and must be refined once the DC line parameters are availabe.
          \item For conductor optimization a lower current rating of the DMR can be applied.
        \end{itemize}
      \end{itemize}
    \end{column}
  \end{columns}
\end{frame}


\begin{frame}{Link overview}
  \centering
  \renewcommand{\arraystretch}{1.3}
  \begin{tabular}{@{}lll@{}}
    \hline
    \textbf{Connection} & \textbf{Topology} & \textbf{Rating} \\
    \hline
    Fennec-Stratos & SMP & 1000 MW \\
    Kit-Corsac & DSMP or BIP & 3000 MW \\
    Marble-Corsac & SMP & 1000 MW \\
    Swift Blanford & DSMP or BIP & 3000 MW \\
    \hline
  \end{tabular}
\end{frame}


% \begin{frame}{Working Packages}
%   \begin{itemize}
%   \setlength{\itemsep}{0.5em}
%   \item \textbf{DC line}
%     \begin{itemize}
%       \item Local company is working on DC route definition and DC line design

%     \end{itemize}
%   \item \textbf{Network Control \& Protection}
%     \begin{itemize}
%       \item Operation concept definition
%       \item Offline studies for concept verification
%       \item Development of models for performance design and verification
%       \item Definition of performance verification tests and evaluation criteria
%       \item Control hardware definition
%     \end{itemize}
%   \item \textbf{HVDC concept evaluation}
%     \begin{itemize}
%       \item Technical evaluation of HVDC concepts
%       \item Commercial evaluation of HVDC concepts
%       \item OEMs should be approached with only one concept
%     \end{itemize}
% \end{itemize}
% \end{frame}
